\documentclass[11pt,a4paper]{article}

\usepackage[utf8]{inputenc}
\usepackage[T1]{fontenc}
\usepackage[brazilian]{babel}
\usepackage{lmodern}
\usepackage[a4paper,margin=2.4cm]{geometry}
\usepackage{xcolor}
\usepackage{hyperref}
\usepackage{enumitem}
\usepackage{listings}
\usepackage{parskip}
\usepackage{needspace}

\hypersetup{colorlinks=true, linkcolor=black, urlcolor=blue!55!black,
  pdftitle={Guia de estudo - Sprint 2}, pdfauthor={Matheus Moreli}}
\setlist{itemsep=3pt, topsep=5pt}
\setlength{\emergencystretch}{3em}
\lstset{basicstyle=\ttfamily\small, backgroundcolor=\color{black!4},
  frame=single, rulecolor=\color{black!20}, breaklines=true,
  columns=fullflexible, showstringspaces=false, xleftmargin=4pt, xrightmargin=4pt}
\newcommand{\projeto}{\paragraph{Aplicação no projeto.}}
\newcommand{\comoaplicar}{\paragraph{Como aplicar.}}
\newcommand{\referencia}[2]{\paragraph{Card e referência.}\href{#1}{#2}}
\newcommand{\cardsection}[1]{\Needspace{12\baselineskip}\subsection{#1}}

\begin{document}

\begin{titlepage}
\thispagestyle{empty}
\centering
\vspace*{3.2cm}
{\Large\textbf{SPRINT 2}}\\[0.5cm]
{\LARGE Inventário de Ativos e Vulnerabilidades}\\[1.2cm]
{\large Material de apoio para apresentação e estudo}\\[3.2cm]
\begin{flushleft}
\large
\textbf{Aluno:} Matheus Moreli Durão Silva\\[0.35cm]
\textbf{Matrícula:} 12621CBS217\\[0.35cm]
\textbf{E-mail:} \href{mailto:matheus.moreli@ufu.br}{matheus.moreli@ufu.br}\\[0.35cm]
\textbf{Curso:} Cibersegurança - Bacharelado - Noturno
\end{flushleft}
\vfill
\end{titlepage}

\tableofcontents
\newpage

\section{O que mudou da Sprint 1 para a Sprint 2}

Na Sprint 1, os conceitos foram praticados em exercícios pequenos: entrada, saída, condições, listas, dicionários, JSON e Git. A Sprint 2 junta essas peças em um programa único. O foco deixa de ser ``como usar um comando isolado'' e passa a ser ``como organizar uma funcionalidade completa sem tornar o código complicado''.

O programa continua sendo de terminal e usa apenas recursos básicos de Python: funções, classes simples, \texttt{Enum}, dicionários, listas, arquivo JSON e \texttt{try/except}. Não há banco de dados, interface gráfica, framework ou consulta automática à internet. Essa escolha é intencional: o problema é pequeno e precisa ser fácil de explicar.

\section{Organização simples do código}

O projeto foi separado por responsabilidade. Separar não significa criar muitas camadas: significa evitar que um mesmo arquivo faça tudo ao mesmo tempo.

\begin{itemize}
  \item \texttt{app.py}: inicia o programa e informa o caminho do JSON.
  \item \texttt{interface.py}: contém \texttt{input()}, \texttt{print()} e o menu.
  \item \texttt{inventario.py}: contém regras de cadastro, consulta, atualização e exclusão.
  \item \texttt{modelos.py}: define \texttt{Ativo}, \texttt{Vulnerabilidade} e \texttt{TipoAtivo}.
  \item \texttt{persistencia.py}: lê e grava o arquivo JSON.
  \item \texttt{test\_sprint2.py}: verifica automaticamente regras importantes.
\end{itemize}

\begin{lstlisting}[language=Python,caption={Fluxo principal simplificado}]
ativos = carregar_ativos("dados/inventario.json")
executar_programa(ativos, "dados/inventario.json")
\end{lstlisting}

\projeto
O menu não precisa conhecer os detalhes de JSON, e a função que salva JSON não precisa fazer perguntas ao usuário. Quando cada parte tem uma tarefa principal, fica mais fácil localizar erros e explicar o código.

\section{Classes e Enum}

\cardsection{S2\_01 e S2\_02 -- Boas práticas, classes e funções}

Uma classe é um molde para representar dados relacionados. Neste trabalho, \texttt{Ativo} reúne ID, nome, responsável, localização, tipo, criticidade, vulnerabilidades e histórico. \texttt{Vulnerabilidade} reúne CVE, CWE, CVSS, descrição, fonte, data, impacto, prioridade, tratamento, status e verificação.

\begin{lstlisting}[language=Python,caption={Ideia de uma classe simples}]
class Ativo:
    def __init__(self, identificador, nome):
        self.id = identificador
        self.nome = nome
        self.vulnerabilidades = []
\end{lstlisting}

\texttt{self} representa o objeto que está sendo criado ou alterado. Assim, dois ativos podem existir ao mesmo tempo, cada um com seus próprios dados. Uma função como \texttt{cadastrar\_ativo()} recebe os dados, verifica regras e cria um objeto. Ela não usa \texttt{input()} nem \texttt{print()}; isso fica na interface.

\texttt{TipoAtivo} é um \texttt{IntEnum}. Ele associa nomes claros a códigos inteiros, como \texttt{SERVIDOR = 2}. Isso atende ao requisito de possuir categorias padronizadas sem espalhar números misteriosos pelo código.

\referencia{https://trello.com/c/ALQSY6p9}{S2\_01 no Trello}; \href{https://trello.com/c/hu3YSI6L}{S2\_02 no Trello}.

\section{CRUD de ativos e vulnerabilidades}

\cardsection{S2\_03, S2\_08 a S2\_11 -- Operações com ativos}

CRUD é uma sigla para quatro operações sobre dados: \emph{Create} (criar), \emph{Read} (consultar), \emph{Update} (atualizar) e \emph{Delete} (excluir). O dicionário \texttt{ativos} usa o ID como chave. Isso permite localizar um ativo diretamente.

\begin{lstlisting}[language=Python,caption={Busca direta pela chave do dicionário}]
def consultar_por_id(ativos, identificador):
    return ativos.get(identificador)
\end{lstlisting}

\texttt{.get()} devolve o valor encontrado ou \texttt{None} quando a chave não existe. Isso evita erro ao buscar um ID ausente. No cadastro, o programa rejeita ID duplicado. Na atualização, o ID não é aceito como campo modificável: nome, responsável, localização, tipo e criticidade podem mudar, mas a identidade do ativo é preservada. Na exclusão, a tela mostra o alvo e exige a palavra \texttt{SIM}; ao remover o ativo, sua lista de vulnerabilidades também desaparece junto.

\cardsection{S2\_12 a S2\_14 -- Vulnerabilidades}

Uma vulnerabilidade é uma fragilidade; ameaça é algo que pode explorar uma fragilidade; risco combina a possibilidade do problema com seu impacto. Por isso, o programa não usa apenas uma nota CVSS: ele também registra impacto no ativo, prioridade, tratamento, status e como a correção será verificada.

Cada vulnerabilidade pertence a um ativo. O par ``ID do ativo + CVE'' a identifica de forma simples: o mesmo CVE pode afetar ativos diferentes, mas não pode aparecer duas vezes no mesmo ativo. As quatro operações também existem para vulnerabilidades: cadastrar, consultar, atualizar e remover.

\begin{lstlisting}[language=Python,caption={Validação simples do CVSS}]
nota = float(cvss)
if not 0.0 <= nota <= 10.0:
    raise DadosInvalidosError("CVSS deve estar entre 0.0 e 10.0.")
\end{lstlisting}

O \texttt{raise} interrompe apenas a operação inválida e a interface mostra uma mensagem compreensível. O menu continua aberto. O sistema também exige uma verificação antes de mudar o status para ``Corrigida'', pois dizer que algo foi corrigido sem evidência não é suficiente.

Como referência de consulta externa, o registro \href{https://nvd.nist.gov/vuln/detail/CVE-2021-44228}{CVE-2021-44228 no NVD} descreve uma vulnerabilidade do Apache Log4j com CVSS 10.0. O programa não consulta o NVD automaticamente: o usuário registra manualmente fonte e data, mantendo a aplicação simples e funcionando sem internet.

\referencia{https://trello.com/c/NdgkJaYg}{S2\_03 no Trello}; \href{https://trello.com/c/zzyFLROU}{S2\_12 no Trello}; \href{https://trello.com/c/xMcbgbvn}{S2\_13 no Trello}; \href{https://trello.com/c/XFF3yzHz}{S2\_14 no Trello}.

\section{Requisitos, validações e critérios de aceite}

\cardsection{S2\_04 a S2\_07 -- Como saber se está certo}

Um requisito diz o que o sistema deve fazer. Um critério de aceite diz como comprovar que ele fez corretamente. Por exemplo, ``cadastrar ativo'' é o requisito; ``o ID precisa ser inteiro, positivo e único, e o ativo deve aparecer na busca'' é o critério verificável.

\begin{itemize}
  \item Campo obrigatório vazio: o programa pede novamente, sem fechar.
  \item ID repetido: o ativo antigo não é alterado e o novo não é cadastrado.
  \item Busca inexistente: aparece ``Ativo não encontrado'', sem erro técnico.
  \item Exclusão cancelada: nada é removido.
  \item JSON inválido: o programa não sobrescreve o arquivo corrompido com uma base vazia.
\end{itemize}

Essas regras aparecem em funções de validação, como \texttt{validar\_texto()}, \texttt{validar\_cve()}, \texttt{validar\_cwe()} e \texttt{validar\_cvss()}. Quando uma regra falha, elas usam \texttt{DadosInvalidosError}; a interface captura essa exceção e apresenta uma mensagem curta.

\referencia{https://trello.com/c/NcSDMJuv}{S2\_04 no Trello}; \href{https://trello.com/c/YSC57Mjy}{S2\_05 no Trello}; \href{https://trello.com/c/5beR1TIT}{S2\_06 no Trello}; \href{https://trello.com/c/Vs045ALu}{S2\_07 no Trello}.

\section{Dicionários, JSON e testes}

\cardsection{S2\_15 -- Dicionários como índice}

O dicionário é a estrutura central do inventário:

\begin{lstlisting}[language=Python,caption={Estrutura principal na memória}]
ativos = {
    10: ativo_servidor,
    20: ativo_notebook
}
\end{lstlisting}

Em vez de percorrer todos os itens para achar o ID 10, o código usa \texttt{ativos.get(10)}. A chave é única, o formato dos objetos é uniforme e as alterações preservam a mesma chave. Ao salvar, os objetos viram dicionários simples; ao carregar, voltam a ser objetos \texttt{Ativo} e \texttt{Vulnerabilidade}.

\cardsection{S2\_18 -- Testes e correções}

\texttt{test\_sprint2.py} não é um segundo programa para o usuário usar. É uma lista de verificações automáticas feita com o módulo padrão \texttt{unittest}. Cada teste cria uma base pequena na memória, executa uma regra e compara o resultado esperado com o obtido.

\begin{lstlisting}[language=Python,caption={Como executar os testes}]
python -m unittest -v test_sprint2.py
\end{lstlisting}

Os 14 testes verificam, entre outros casos, ID duplicado, tipo inválido, filtro sem alteração da base, atualização preservando ID, confirmação de exclusão, CVE duplicado, CVSS inválido, CRUD de vulnerabilidades e proteção contra JSON corrompido. Eles não usam dados reais, internet ou senhas.

\section{GitHub e roteiro da apresentação}

\cardsection{S2\_16, S2\_17 e S2\_19 -- Entrega}

O código está no GitHub, na branch \texttt{main}. O histórico possui branches de funcionalidade, commits com mensagens explicativas e merges por Pull Request. A tag \texttt{trabalho-1-entrega} identifica a versão preparada para apresentação.

Para demonstrar em cinco minutos:
\begin{enumerate}
  \item Execute \texttt{python app.py}.
  \item Cadastre um ativo e uma vulnerabilidade associada.
  \item Consulte o ativo e mostre que a vulnerabilidade está vinculada a ele.
  \item Atualize o responsável ou o status da vulnerabilidade.
  \item Tente repetir um ID ou informar CVSS inválido para mostrar a validação.
  \item Remova uma vulnerabilidade, primeiro cancelando e depois confirmando.
  \item Abra o GitHub, mostre \texttt{main}, as Pull Requests e a tag final.
\end{enumerate}

Ao explicar, prefira frases concretas: ``usei um dicionário porque o ID é único e a busca fica direta''; ``usei JSON porque os dados continuam depois que o programa fecha''; ``separei interface e regras para não misturar perguntas com validação''.

\referencia{https://trello.com/c/UvtQ6Mc6}{S2\_16 no Trello}; \href{https://trello.com/c/FBSYXNJP}{S2\_17 no Trello}; \href{https://trello.com/c/ADSVtVQ4}{S2\_19 no Trello}.

\section{Referências gerais}
\begin{itemize}
  \item Cards S2\_01 a S2\_19 do quadro Trello da disciplina.
  \item \href{https://docs.python.org/3/}{Documentação oficial do Python}.
  \item \href{https://docs.python.org/3/library/unittest.html}{Documentação de unittest}.
  \item \href{https://nvd.nist.gov/vuln/detail/CVE-2021-44228}{NVD - CVE-2021-44228}.
  \item \href{https://git-scm.com/book/pt-br/v2}{Pro Git}.
\end{itemize}

\end{document}
